\color{uniblue}\subsection[Выбор уставок защиты от обрыва провода]{Выбор уставок защиты от обрыва провода}
\color{black}

\begin{enumerate}

\item
Принимаем параметр $K_{\textsl{несим}} = 0,35$ в соответствии с пунктом \ref{zop:punkt}. Контроль абсолютного значения тока $I_{2}$ не используется.

\item
Выдержка времени выбирается на ступень селективности больше времени срабатывания самой медленно действующей защиты трансформатора по выражению (\ref{eq:zop}):
\begin{equation*}
    t_{\textsl{с.з.ЗОП}} = t_{\textsl{с.з.МТЗ}}+{\Delta t} = 1,9 + 0,3 = 2,2\; (\textsl{с}),
\end{equation*}

\item
Параметр \textbf{$I_{\textsl{ср}}$} <<Ток срабатывания>> функции ЗОП принимается \textbf{равным 0,35}.
\item
Параметр \textbf{$T_{\textsl{ср}}$} <<Выдержка времени срабатывания>> функции ЗОП принимается \textbf{равным 2,2 c}.
\end{enumerate}